%======================================================================
\section{The LQD family}
\label{sec:model}

\subsection{Definition}

It remains to choose the smooth part $g$, and the choice is guided by two
requirements: the endpoint values must be readable --- they will turn out
to be the tail scales --- and the body must be flexible, well-conditioned,
and cheap.  A polynomial in the rank, expressed in the Legendre basis on
$u\in[0,1]$, meets all three.  The central definition of the paper is the
following; it is the only boxed display.

\begin{definition}[The LQD slice]
\label{def:lqd}
Fix an order $N\ge2$ and a coefficient vector
$\theta=(L,R,a_2,\dots,a_N)\in\R^{N+1}$.  Let $Z$ be a standard logistic
variable, with CDF $\Lambda(z)=1/(1+e^{-z})$ and density
$\rho(z)=\Lambda(z)\Lambda(-z)$.  The LQD slice is the law of
\begin{equation}
\boxed{\;
X=x(Z),\qquad
x(z)=m+\int_{0}^{z}e^{\,g(\Lambda(t))}\,\dd t,\qquad
g(u)=(1-u)L+uR+\sum_{n=2}^{N}a_n P_n(1-2u),
\;}
\label{eq:lqd}
\end{equation}
where $P_n$ is the Legendre polynomial of degree $n$ and the scalar $m$ is
fixed by the martingale condition $\E[e^{X}]=1$.
\end{definition}

Three remarks unpack the display.  First, the map $x$ is strictly
increasing for \emph{every} $\theta$, since its derivative is an
exponential; so $X$ always has a genuine law, with CDF and density
\begin{equation}
F_X\big(x(z)\big)=\Lambda(z),
\qquad
f_X\big(x(z)\big)=\frac{\rho(z)}{x'(z)}
=\Lambda(z)\Lambda(-z)\,e^{-g(\Lambda(z))}\;>\;0 .
\label{eq:density}
\end{equation}
This single line is the entire no-arbitrage mechanism of
\cref{sec:validity}; no positivity condition is ever checked, because none
can fail.  Second, the quantile function of $X$ is $Q(u)=x(\logit u)$, and
by \eqref{eq:speed} its quantile density is
\begin{equation}
q(u)=\frac{e^{g(u)}}{u(1-u)},
\qquad\text{i.e.}\qquad
\log q(u)=-\log u-\log(1-u)+g(u):
\label{eq:lqdform}
\end{equation}
the definition is exactly the skeleton-plus-smooth-part decomposition
\eqref{eq:skeleton} with a polynomial $g$ --- hence the name
\emph{log-quantile-density} model.  The transport form \eqref{eq:lqd} and
the LQD form \eqref{eq:lqdform} are the same object; the paper computes in
the first and reasons about distributions in the second.  Third, the
expansion starts at $n=2$ deliberately: degrees zero and one are already
spanned by the \emph{endpoint chord} $(1-u)L+uR$, since
$(1-u)L+uR=\frac{L+R}{2}+\frac{L-R}{2}(1-2u)$ and $P_0=1$, $P_1(1-2u)=1-2u$.
The family at order $N$ is thus exactly the set of laws whose $g$ is a
polynomial of degree $\le N$; the chord is a choice of \emph{coordinates}
on that space, made so that the two leading coefficients sit at the
endpoints, where the tails live.  Nothing in the definition needs $N$
larger than $2$; the reference implementation restricts to $N\ge4$, a
calibration convention aligned with the ridge of \cref{sec:calibration}
(which starts penalizing at $n=4$) rather than a mathematical
requirement.  Legendre itself is also a disclosed convention: any
polynomial basis spans the same family, but the Legendre modes are
orthogonal under the uniform rank measure --- the inner product every
rank integral in the model uses --- so high-order coordinates remain
numerically distinguishable in the least-squares normal equations,
where monomials in $u$ would reproduce Hilbert-matrix conditioning.
Chebyshev would serve comparably; monomials would not.

Because $P_n(1)=1$ and $P_n(-1)=(-1)^{n}$, the two endpoint values of $g$
are
\begin{equation}
g(0)=L+\sum_{n\ge2}a_n,
\qquad
g(1)=R+\sum_{n\ge2}(-1)^{n}a_n ,
\label{eq:endpoints}
\end{equation}
and we name their exponentials the \emph{endpoint speeds}
\begin{equation}
\lambda_-=e^{g(0)},
\qquad
\lambda_+=e^{g(1)} .
\label{eq:speeds}
\end{equation}
By \eqref{eq:speed} these are the asymptotic speeds of the transport.
Precisely, writing $\bar{x}(z)=\int_0^z e^{g(\Lambda(t))}\dd t$ for the
\emph{unshifted} transport (so $x=m+\bar{x}$), there are constants
$b_-$, $b_+$ with
\begin{equation}
\bar{x}(z)=b_-+\lambda_-z+O(e^{z})\ \ (z\to-\infty),
\qquad
\bar{x}(z)=b_++\lambda_+z+O(e^{-z})\ \ (z\to+\infty),
\label{eq:asymptotes}
\end{equation}
because $g$ is Lipschitz on $[0,1]$ and $\Lambda(z)$ approaches its
endpoint at the rate $e^{-|z|}$, so the speed $e^{g(\Lambda(z))}$
approaches $\lambda_\pm$ at that same rate and the deviation integrates
to a convergent constant plus an exponentially small tail.  The shifted
transport obeys the same asymptotes with constants $m+b_\pm$; keeping
the two families of constants apart matters below (the tail prefactor
of \cref{prop:tails} contains $m+b_+$, not $b_+$).  This estimate does
real work in \cref{prop:shift,sec:tails}.  Note, for later honesty, that
\eqref{eq:endpoints} couples body and tails in the raw coordinates: every
$a_n$ moves one or both endpoint values.  Adding $0.10$ to $a_2$ multiplies
\emph{both} speeds by $e^{0.10}\approx1.105$.  The fitted law is
coordinate-independent, but the raw coordinates are a poor language for
tail priors; \cref{sec:computation} introduces the two reparametrizations
--- the endpoint chart and the logistic chart --- that repair this.

What does one coefficient do?  Raising $g$ near a rank $u$ stretches the
return spacing there, and stretched spacing is thinned probability
\eqref{eq:reciprocal}: a basis mode is a mass-shuffling instruction written
in rank coordinates.  Even modes ($P_2$, $P_4$, \dots) act symmetrically
--- $a_2>0$ stretches both tails and compresses the center, a sharper peak
with fatter tails, which the smile shows as extra convexity; odd modes tilt
--- $a_3>0$ slides mass from one side to the other, which the smile shows
as skew; higher modes sculpt with finer brushes.  \Cref{fig:modes} shows
the three charts of the same three perturbations.

\begin{figure}[htbp]
\centering
\paperfig{\textwidth}{fig_modes.pdf}
\caption{One coefficient at a time: $a_2$, $a_3$, or $a_4$ set to $0.10$ on
a symmetric constant-speed reference, everything else untouched.  Panel A:
the perturbation of the log-speed $g$ over the rank --- even modes
symmetric, odd modes not.  Panel B: the implied-volatility smile --- mode 2
bends it symmetrically, mode 3 tilts it, mode 4 works the shoulders.
Panel C: the density --- mass thins where $g$ rises and thickens where it
falls, and every perturbed density is non-negative without anything being
checked.  Every raw mode drawn here also moves \emph{both} tail scales, by
a factor $e^{\pm0.10}$: the even modes multiply both by $e^{0.10}$, while
$a_3$ multiplies $\lambda_-$ by $e^{0.10}$ and $\lambda_+$ by $e^{-0.10}$
--- $P_n(-1)=(-1)^{n}$ in \eqref{eq:endpoints}, visible in panel A as the
$a_3$ curve ending below baseline at $u=1$.  This is the coupling the
endpoint chart of \cref{sec:computation} removes.}
\label{fig:modes}
\end{figure}

\subsection{The martingale shift}

The scalar $m$ deserves one short proposition, because it is the only
normalization in the model and everything else stands on it.  Recall
$\bar{x}$, the unshifted transport of \eqref{eq:asymptotes}, with
$x=m+\bar{x}$.

\begin{proposition}[Martingale normalization]
\label{prop:shift}
Let $I=\int_{-\infty}^{\infty}e^{\bar{x}(z)}\rho(z)\,\dd z$.  If
$I<\infty$, then $m=-\log I$ is the unique scalar for which
$\E[e^{X}]=1$, and the shift changes neither the density's shape, nor its
tails, nor its positivity.  Moreover $I<\infty$ if and only if
$\lambda_+<1$.
\end{proposition}

\begin{proof}
Since $X=x(Z)$ and $Z$ has density $\rho$,
$\E[e^{X}]=\int e^{m+\bar{x}(z)}\rho(z)\,\dd z=e^{m}I$; if $I$ is finite,
$e^{m}I=1$ has the unique solution $m=-\log I$.  The shift translates the
law of $X$ by a constant, which multiplies $f_X$ by nothing and moves every
quantile equally.  For the convergence claim, use \eqref{eq:asymptotes}:
as $z\to+\infty$, $\bar{x}(z)=b_++\lambda_+z+O(e^{-z})$ with $b_+$
finite, so the integrand is asymptotically $e^{b_+}$ times
$e^{(\lambda_+-1)z}$, using $\rho(z)\sim e^{-z}$.  This is integrable at
$+\infty$ iff $\lambda_+<1$; if $\lambda_+=1$ the integrand tends to a
positive constant, and if $\lambda_+>1$ it grows, so the integral diverges
in either case.  As $z\to-\infty$, symmetrically,
$\bar{x}(z)=b_-+\lambda_-z+O(e^{z})$ with $\lambda_->0$, so the integrand
is asymptotically a constant times $e^{(1+\lambda_-)z}$, integrable at
$-\infty$ for every $\lambda_->0$: the left end never obstructs.
\end{proof}

The condition $\lambda_+<1$ is the \emph{finite-forward wall}, and it is
the single validity constraint of the whole construction;
\cref{sec:validity} discusses its meaning and \cref{sec:tails} shows it is
not an artifact of the basis but a fact about exponential tails.

\subsection{An exactly solvable slice}
\label{sec:constspeed}

Set every body coefficient to zero and the two chord coefficients equal:
$a_n=0$, $L=R=\log s$ with $0<s<1$.  Then $g\equiv\log s$, the transport is
exactly affine, $x(z)=m+sz$, and $X$ is a logistic random variable with
scale $s$.  This slice is to the LQD family what the harmonic oscillator is
to mechanics: every formula of the paper can be evaluated on it in closed
form, and \cref{sec:examples} uses it as a machine-precision audit of the
implementation.

Everything flows from one integral.  The moment generating function of the
standard logistic score is a rank integral that Euler solved: substituting
$u=\Lambda(z)$, $e^{tz}=\big(\frac{u}{1-u}\big)^{t}$,
\begin{equation}
\E\big[e^{tZ}\big]
=\int_0^1\Big(\frac{u}{1-u}\Big)^{t}\dd u
=\mathrm{B}(1+t,\,1-t)
=\Gamma(1+t)\Gamma(1-t)
=\frac{\pi t}{\sin(\pi t)},
\qquad |t|<1,
\label{eq:mgf}
\end{equation}
where $\mathrm{B}$ is the beta function, the middle equality is its Gamma
representation, and the last is Euler's reflection formula.  Note that
\eqref{eq:mgf} already confirms the wall: the MGF exists exactly on the
strip $|t|<1$, as \cref{sec:tails} says it must for a law with unit
endpoint speeds.  Three consequences:
\begin{enumerate}[leftmargin=1.6em]
\item \emph{The shift in closed form.}  $I=\E[e^{sZ}]=\pi s/\sin(\pi s)$,
so $m=-\log\big(\pi s/\sin(\pi s)\big)$ --- the only slice in the family
whose normalization needs no quadrature, and therefore the natural
end-to-end test vector.
\item \emph{The variance, hence the initializer.}  Expanding
$\log\E[e^{tZ}]=\frac{\pi^2t^2}{6}+O(t^4)$ (from
$\sin\pi t=\pi t(1-\pi^2t^2/6+\cdots)$), the second cumulant is
$\operatorname{Var}(Z)=\pi^{2}/3$, so
$\operatorname{Var}(X)=\pi^{2}s^{2}/3$.  Matching a target ATM total
variance $w_0$ gives the cold-start seed
\begin{equation}
s=\frac{\sqrt{3w_0}}{\pi},
\label{eq:coldstart}
\end{equation}
which is where every calibration in this paper starts when it has no better
information (\cref{sec:calibration}).
\item \emph{The seed's known bias.}  Matching the \emph{variance} does not
match the ATM \emph{price}.  For small $s$, the at-the-money call of this
slice is $c(0)=\E[(e^{X}-1)^{+}]\approx s\,\E[Z^{+}]=s\log2$ (two small
steps deserve their receipts: the shift is negligible at this order,
$m=-\log\frac{\pi s}{\sin\pi s}=-\frac{\pi^{2}s^{2}}{6}+O(s^{4})$, so
$(e^{X}-1)^{+}\approx X^{+}\approx sZ^{+}$; and
$\E[Z^{+}]=\int_0^\infty z\,\rho(z)\,\dd z
=\sum_{n\ge1}\frac{(-1)^{n-1}}{n}=\log 2$, by expanding
$\rho(z)=e^{-z}(1+e^{-z})^{-2}$ in a geometric series and integrating
term by term), while the Black price at total variance $w$ is
$\approx\sqrt{w/2\pi}$; equating gives an implied ATM total variance
$w\approx2\pi(\log2)^{2}s^{2}$, and the ratio to the variance-matched
target is
\[
\frac{2\pi(\log 2)^{2}}{\pi^{2}/3}
=\frac{6(\log 2)^{2}}{\pi}\approx0.918 .
\]
The cold start therefore under-prices ATM variance by about $8\%$ --- a
derived constant, confirmed in panel B of \cref{fig:exact} --- which the
first optimizer iteration removes.  It is recorded here because a seed
whose bias is known is a seed that can be trusted.
\end{enumerate}

Admissibility of the slice is $\lambda_-=\lambda_+=s<1$; a $55\%$
volatility at one year is $s\approx0.30$, nowhere near the wall.
\Cref{fig:exact} confronts the reference implementation with the closed
forms across scales $s\in[0.05,0.95]$: the transport map agrees with
$m+sz$ to \MacExactMapWorst\ (cumulative-Simpson roundoff), the shift
with $-\log(\pi s/\sin\pi s)$ to \MacExactMuWorst, and the martingale
check $|\E[e^{X}]-1|$ stays below \MacExactMartWorst.  The three levels
are worth distinguishing rather than blurring into ``machine
precision'': the shift error is the largest and \emph{grows toward the
wall}, as it must --- the normalizing integrand decays like
$e^{(\lambda_+-1)z}$, so as $s\to1$ its tail thickens and the
truncation-corrected quadrature works harder.  A wrong grid constant, a
mis-anchored cumulative sum, or a stale cache moves any of these
agreements by many orders of magnitude, which is what makes the audit
sharp.

\begin{figure}[htbp]
\centering
\paperfig{\textwidth}{fig_exact.pdf}
\caption{The constant-speed audit, in two panels.  Panel A: the
implementation's martingale shift $m$ and transport map $x(z)$ against
the closed forms $-\log(\pi s/\sin\pi s)$ and $m+sz$ from
\eqref{eq:mgf}, absolute error on a log scale across the admissible
scale range: map error at most \MacExactMapWorst; shift error at most
\MacExactMuWorst, climbing as $s\to1$ because the normalizing
integrand's decay rate $1-\lambda_+$ vanishes at the wall.  Panel B:
the cold-start bias of \eqref{eq:coldstart} --- the ATM variance
mismatch of the variance-matched seed, measured
\MacColdStartMismatchPct\% at the $20\%$-vol slice, against the derived
small-$s$ limit $6(\log2)^2/\pi-1\approx-8.2\%$; the first calibration
step removes it.}
\label{fig:exact}
\end{figure}

\subsection{A fitted slice in the model's own chart}

Before proving anything, it is worth seeing what a real fitted $g$ looks
like.  \Cref{fig:lqdchart} takes one liquid equity-index expiry from the
frozen snapshot of \cref{sec:examples} (SPY, December 2026; provenance in
\cref{app:data}) and draws the fitted $\log q(u)$: the universal skeleton
$-\log u-\log(1-u)$ dashed, the fitted smooth part $g$ shaded between the
curves.  The entire freedom of the model is that shaded band ---
low-degree, gently sloped, its two endpoint values the fitted tail scales
--- and \emph{any} smooth band whatever, shifted, bent, or oscillating, is
another arbitrage-free slice provided only that the right-hand endpoint
stays below zero.  The companion panel shows the same object as a density.

\begin{figure}[htbp]
\centering
\paperfig{\textwidth}{fig_lqd_chart.pdf}
\caption{A real fitted slice (SPY, December 2026 expiry, from the frozen
snapshot of \cref{app:data}) in the model's own chart.  Panel A: fitted
$\log q(u)$ (solid) over the universal skeleton $-\log u-\log(1-u)$
(dashed); the shaded gap between them is the fitted $g$ --- the model's
entire freedom.  Its left endpoint sits higher than its right: a fatter
left tail, the index put skew in the model's native coordinates.  Panel B:
the same slice as a return density, skewed left, with the strikes of the
quoted strip marked --- the region where the data, rather than the basis,
chose the curve.}
\label{fig:lqdchart}
\end{figure}
